\section{Теорема о невозможности тепловой смерти \texorpdfstring{$\Mlayer$}{M}}
\label{sec:theorem-23}

\begin{definition}[M-эволюция]
\label{def:state-evolution}
$\symCapitalPsi^{t+1}=g(\symCapitalPsi^t)$ на счётном $\carrierLattice$; один $\htick$ на тик; $g$ постоянен.
\end{definition}

\begin{definition}[Останов · D1]
\label{def:D1}
Существует $t$ без шага $g$, или $\htick\to\infty$, или $g$ меняется / обнуляется.
\end{definition}

\begin{definition}[M heat death · D2]
\label{def:D2}
$\sum_x|z|^2\to 0$, или глобально $z\equiv 0$ устойчиво, или $g(\symCapitalPsi^*)=\symCapitalPsi^*$ с «мёртвым'' вакуумом без A5-кипения.
\end{definition}

\begin{definition}[$\Tlayer$ heat death · D3]
\label{def:D3}
Макроосреднение $\Phi=\mathcal{B}z$ теряет структуру ($\mathrm{Var}(\Phi)\to 0$) при $N\approx\mathrm{const}$ на $\Mlayer$.
\end{definition}

\begin{lemma}[Нет останов]
\label{lem:231}
В постановке M-эволюции (\cref{def:state-evolution}) для всех $t\in\mathbb{Z}$: $\symCapitalPsi^{t+1}=g(\symCapitalPsi^t)$, $\htick$ и $g$ постоянны.
\end{lemma}

\begin{proof}
\pusing{определения~\ref{def:state-evolution}}
\[
  \symCapitalPsi^{t+1}=g(\symCapitalPsi^t)\quad\text{для всех }t\in\mathbb{Z},
\]
при постоянных $\htick$ и~$g$.
\pfrom{определения D1 (\ref{def:D1})}
остановка потребовала бы прекращения~$g$, $\htick\to\infty$ или смены закона ---
это уже другая постановка, не M-эволюции (\cref{def:state-evolution}).
\end{proof}

\begin{lemma}[Сохранение информации]
\label{lem:232}
$N(t+1)=N(t)$ для всех $t$.
\end{lemma}

\begin{proof}
\pusing{определения поля (\cref{def:field-a3})}
\[
  N(t)=\sum_x z^\dagger(x,t)z(x,t).
\]
Линейная часть~$g$ --- произведение локальных унитарных операторов на рёбрах~$N(x)$;
каждый фактор сохраняет $z^\dagger z$ на затронутых узлах.
Нелинейная часть --- только $z\mapsto z\,e^{i\Phi(z)}$, также сохраняющая $|z|$.
\pfollows
композиция шага не меняет $\sum_x|z|^2$, то есть $N(t+1)=N(t)$.
\end{proof}

\begin{corollary}
\label{cor:232a}
$N\to 0$ на $\Mlayer$ невозможно, если $N(0)>0$.
\end{corollary}

\begin{proof}
\pfrom{леммы~\ref{lem:232} при $N(0)>0$}
\[
  N(t)=N(0)>0\qquad\text{для всех }t\in\mathbb{Z}.
\]
\phence предел $N\to 0$ недостижим.
\end{proof}

\begin{lemma}[Фазовая энтропия и норма]
\label{lem:233}
Рост фазовой энтропии не влечёт $N\to 0$.
\end{lemma}

\begin{proof}
\pusing{фиксированной нормы $N$}
меняются только аргументы~$z$; $|z|$ сохраняется.
\pfollows
рост фазовой энтропии не даёт механизма $N\to 0$.
\end{proof}

\begin{lemma}[Нет устойчивого $z\equiv 0$]
\label{lem:234}
$z(x,t)\equiv 0$ не устойчивое состояние $g$.
\end{lemma}

\begin{proof}
\pfrom{тождества $z\equiv 0$}
фаза $z e^{i\Phi}$ остаётся нулём при любой~$\Phi$:
локальный шаг не выводит поле из «мёртвой'' точки.
\pusing{\Axref{ax:A5}}
вакуум --- осциллирующий фон с $\rho_{\mathrm{field}}>0$.
\phence глобально устойчивый $z\equiv 0$ исключён; вместе с \cref{cor:232a} ветвь $N\to 0$ D2 закрыта.
\end{proof}

\begin{lemma}[Микрообратимость]
\label{lem:235}
$g$ обратима на $Z_N[i]$.
\end{lemma}

\begin{proof}
\pfrom{\Axref{ax:A13}}
существует $g^{-1}$ с $g^{-1}(g(\symCapitalPsi))=\symCapitalPsi$ на $\mathbb{Z}_{N_{\mathrm{ring}}}[i]$.
\pusing{схемы «лягушка'' (\cref{sec:congruence}, CL-1)}
\[
  Z^+(x,t)+Z^-(x,t)=2Z(x,t)+\lfloor\mathcal{N}(x,t)\rfloor,
\]
откуда явно восстанавливается $Z^-$ при учёте фазовых импульсов~$\Phi$.
\pfollows
микрошаг --- перестановка на конечном $\mathcal{C}$.
\end{proof}

\begin{lemma}[$\mathcal{B}$ необратима]
\label{lem:236}
При $R\ge 1$ осреднение $\mathcal{B}$ не инъективен.
\end{lemma}

\begin{proof}
\puseref{осреднения с биномиальным ядром}{def:soft-readout}
\[
  \Phi=\mathcal{B}z,\qquad W=\tfrac14[1,2,1]^{\otimes 2R}
\]
(на срезе) или эквивалентное $3\times 3$-ядро.
\pfrom{$W\ne\delta$ и конечности $\carrierLattice$}
существуют $z\ne z'$ с $\mathcal{B}z=\mathcal{B}z'$.
\phence $\mathcal{B}$ не инъективна при $R\ge 1$.
\end{proof}

\begin{proposition}[D3 не нарушает $\Mlayer$]
\label{prop:237}
$\nu_{\mathrm{CA}}$-затухание на $\Tlayer$ --- не член $g$; $N$ invariant (лемма~\ref{lem:232}).
\end{proposition}

\begin{proof}
\emph{Первое утверждение.}
\pusing{структуры~$g$ на~$\Mlayer$}
закон --- композиция локальных унитарных преобразований
и фаз $e^{i\Phi(z)}$; в нём нет вещественного члена $\nu_{\mathrm{CA}}\nabla^2 z$,
 иначе $|z|$ перестал бы быть инвариантом в нарушение \Axref{ax:A3}.
Кинематическая вязкость $\nu_{\mathrm{CA}}\nabla^2\Phi$ появляется только в эффективном
уравнении для $\Phi=\mathcal{B}z$ на~$\Tlayer$:
это остаток дискретного осреднения при конечных $\caStepSpace$, $\htick$.

\smallskip
\emph{Второе утверждение.}
Макроскопическое «затухание'' --- потеря различимости после необратимого~$\mathcal{B}$
(лемма~\ref{lem:236}): прибор не видит часть микрофазы.
На~$\Mlayer$ эта энергия не исчезает из~$N(t)$ (лемма~\ref{lem:232}):
по \Axref{ax:A5} она уходит в микро-«кипение'' --- перераспределение фаз
при $|z|\ge z_{\min}>0$; по \Axref{ax:A6} --- в рост локальной фазовой энтропии
при фиксированном~$N$.
Сценарий D3 ($\mathrm{Var}(\Phi)\to 0$ при $N\approx\mathrm{const}$) описывает
грубый readout на~$\Tlayer$.
\end{proof}

\begin{definition}[Frozen tick · D4]
\label{def:D4}
$\lfloor\mathcal{N}(\symCapitalPsi)\rfloor\equiv 0$ на всём $\carrierLattice$.
\end{definition}

\begin{definition}[Глобальная неподвижная точка · D5]
\label{def:D5}
$g(\symCapitalPsi^*)=\symCapitalPsi^*$ при $(Z_t,Z_{t-1})=(Z^*,Z^*)$ $\Leftrightarrow$ D4.
\end{definition}

\begin{lemma}[Классификация D5]
\label{lem:238a}
D5 $\Leftrightarrow$ $z(x,t)\equiv c$ на торе $\carrierLattice$, $c\in Z_N[i]$.
\end{lemma}

\begin{proof}
\emph{($\Rightarrow$).}
\pfrom{D5 (\ref{def:D5})}
$\lfloor\mathcal{N}\rfloor\equiv 0$ на всём $\carrierLattice$.
\pfollows
$\Phi=0$ на столкновениях; по \Axref{ax:A9} и связности~$N$
\[
  z(x)\equiv c\in Z_N[i].
\]

\smallskip
\emph{($\Leftarrow$).}
\pfrom{$z(x)\equiv c$}
все $\lfloor\mathcal{N}\rfloor$ обнуляются, $(Z_t,Z_{t-1})$ не меняются.
\phence $g(\symCapitalPsi^*)=\symCapitalPsi^*$.
\end{proof}

\begin{lemma}[A5/A6 исключают физические D5]
\label{lem:238b}
Ни одна конфигурация из леммы~\ref{lem:238a} не A5-физична.
\end{lemma}

\begin{proof}
\pusing{классификации леммы~\ref{lem:238a}}
пусть $z\equiv c$ --- неподвижная точка.
При $c=0$ --- противоречие \Axref{ax:A5}.
При $c\ne 0$ поле статично на всех тактах, что нарушает \Axref{ax:A5} и \Axref{ax:A6}.
\phence ни один случай не A5-физичен.
\end{proof}

\begin{theorem}[Нет физических неподвижных точек без кипения]
\label{thm:238}
Класс A5-физических глобальные неподвижные точки $g(\symCapitalPsi^*)=\symCapitalPsi^*$ без микрокипения --- пуст.
\end{theorem}

\begin{proof}
Пусть $g(\symCapitalPsi^*)=\symCapitalPsi^*$ --- A5-физичная точка без микрокипения.
\pfrom{леммы~\ref{lem:238a}}
$z(x)\equiv c$ на связном~$\carrierLattice$.
\pfrom{леммы~\ref{lem:238b}}
ни $c=0$, ни статичное $c\ne 0$ не допустимы.
\phence таких точек нет.
\end{proof}

\begin{figure}[htbp]
  \centering
  \begin{tikzpicture}[carrier panel, x=0.95cm, y=2.2cm]
  \draw[carrier object, domain=0:10, samples=80] plot (\x, {1});
  \draw[carrier muted, dashed, domain=0:10, samples=80]
    plot (\x, {0.8*exp(-0.25*\x) + 0.05*sin(deg(2.5*\x)) + 0.08});
  \node[anchor=west, font=\scriptsize] at (0.2,1.05) {$N(t)$ на $\mathcal{M}$};
  \node[anchor=west, carrier muted, font=\scriptsize] at (6.5,0.35) {$\mathrm{Var}(\symCapitalPhi)$ на $\mathcal{T}$};
  \node[carrier muted, font=\scriptsize, rotate=90] at (-0.55,0.55) {нормированная величина};
  \node[carrier muted, font=\scriptsize] at (5,-0.25) {время $t$};
  \node[anchor=north, font=\footnotesize] at (5,-0.55)
    {$g^{-1}$ существует; $\opBoundary$ необратима};
\end{tikzpicture}

  \caption{Инвариант $N(t)$ на $\Mlayer$ и возможное затухание контраста $\mathrm{Var}(\Phi)$ на $\Tlayer$.}
  \label{fig:axiom-heat-death}
\end{figure}

\begin{theorem}[Тепловая смерть $\Mlayer$ невозможна]
\label{thm:main-23}
Пусть на~$\Mlayer$ норма сохраняется, вакуум осциллирует,
фазы перемешиваются при фиксированной норме, и у локального закона~$g$ есть обратный шаг.
Тогда:
\begin{enumerate}[label=(\roman*)]
\item \emph{Нет останов} (D1): $\symCapitalPsi^{t+1}=g(\symCapitalPsi^t)$ на каждом такте при постоянных $\htick$ и~$g$.
\item \emph{Нет тепловой смерти на~$\Mlayer$} (D2): не выполняется ни $N\to 0$,
      ни глобально устойчивое $z\equiv 0$,
      ни «мёртвая'' неподвижная точка $g(\symCapitalPsi^*)=\symCapitalPsi^*$ без микро-кипения.
\item \emph{Фазовая энтропия и норма}: рост фазовой энтропии не даёт $N\to 0$.
\item \emph{Микрообратимость, макронеобратимость}: $g$ обратим на решётке;
      осреднение~$\mathcal{B}$ необратимо; кинематическая вязкость на~$\Tlayer$ не входит в~$g$
      и не обнуляет~$N$ на~$\Mlayer$.
\item \emph{Нет A5-физических D5}: глобальные неподвижные точки «лягушки'' не физичны.
\end{enumerate}
\end{theorem}

\begin{proof}
Работаем в постановке M-эволюции (\cref{def:state-evolution}): один такт~$\htick$, постоянный локальный закон~$g$.

\smallskip
\emph{(i) $\neg$D1.}
Из постановки M-эволюции (\cref{def:state-evolution}) следует для каждого $t\in\mathbb{Z}$ выполняется $\symCapitalPsi^{t+1}=g(\symCapitalPsi^t)$,
шаг~$\htick$ и правило~$g$ не зависят от~$t$.
Ни прекращение шага, ни $\htick\to\infty$, ни смена/обнуление~$g$ не следуют из такой эволуции ---
это было бы изменением модели (лемма «Нет останов'', \ref{lem:231}).

\smallskip
\emph{(ii) $\neg$D2.}
Определение~\ref{def:D2} перечисляет три ветви «тепловой смерти'' на~$\Mlayer$.
Первая, $\sum_x|z|^2=N\to 0$, противоречит сохранению нормы:
при $N(0)>0$ лемма «Сохранение информации'' (\ref{lem:232}) даёт $N(t)=N(0)>0$ на всех тактах
(следствие, \ref{cor:232a}), так что предел $N\to 0$ недостижим.
Вторая ветвь --- глобально устойчивый $z\equiv 0$.
При нулевом поле фазовый поворот $z e^{i\Phi}$ остаётся нулём, эволюция не выводит конфигурацию
из этой точки; $|z|=0$ --- вырожденная точка, вакуум осциллирует
(лемма «Нет устойчивого $z\equiv 0$'', \ref{lem:234}).
Третья ветвь --- неподвижная точка $g(\symCapitalPsi^*)=\symCapitalPsi^*$ с «мёртвым'' вакуумом без микро-кипения.
Любая такая точка «лягушки'' есть однородное поле $z(x)\equiv c$ (лемма «Классификация D5'', \ref{lem:238a});
\ref{lem:238b} исключает $c=0$ (мёртвый вакуум) и статичное $c\ne 0$
(нет осциллирующего фона, и фазы глобально заморожены).
Класс A5-физичных неподвижных точек без кипения пуст (теорема «Нет физических неподвижных точек без кипения'', \ref{thm:238}).
Все три ветви D2 закрыты.

\smallskip
\emph{(iii) Фазовая энтропия и норма.}
При фиксированной норме~$N$ перераспределяются фазы;
модуль $|z|$ сохраняется.
Рост фазовой энтропии не содержит механизма $N\to 0$ (лемма «Фазовая энтропия и норма'', \ref{lem:233}).

\smallskip
\emph{(iv) Микрообратимость и макронеобратимость.}
На~$\Mlayer$ шаг~$g$ обратим: существует $g^{-1}$ с $g^{-1}(g(\symCapitalPsi))=\symCapitalPsi$
на $\mathbb{Z}_{N_{\mathrm{ring}}}[i]$ (лемма «Микрообратимость'', \ref{lem:235}).
Осреднение $\Phi=\mathcal{B}z$ при $R\ge 1$ не инъективно --- разные микросостояния дают одно макроосреднение
(лемма «$\mathcal{B}$ необратима'', \ref{lem:236}).
Кинематическая вязкость $\nu_{\mathrm{CA}}\nabla^2\Phi$ входит только в эффективное уравнение на~$\Tlayer$
(предложение «D3 не нарушает $\Mlayer$'', \ref{prop:237}). Макроскопическое «затухание'' --- потеря различимости после~$\mathcal{B}$,
тогда как $N(t)$ на~$\Mlayer$ сохраняется и уходит в микро-перераспределение фаз.

\smallskip
\emph{(v) Нет A5-физических D5.}
Неподвижная точка D5 эквивалентна D4 (\ref{def:D5}) и сводится к $z\equiv c$ (\ref{lem:238a});
ни одна такая конфигурация не осциллирует и не перемешивает фазы одновременно (\ref{lem:238b}),
откуда \ref{thm:238}.
\end{proof}
